\subsection{Reporting prediction retains value under geographic holdout}
All eleven conventional configurations were evaluated on the same 30 spatial test folds (Table~\ref{tab:spatialvalid}). The published GLM m4 had mean AUC 0.800, AP 0.253 and Brier score 0.0401. Terrain-augmented LightGBM had the highest mean AUC, 0.846. Random forest had the highest mean AP, 0.286, and lowest mean Brier score, 0.0384. Its top-ranked 10\% of test checklists contained 53.4\% of positive reports, compared with 46.1\% for m4. Relative to m4, random forest increased AP by approximately 13.1\%, reduced Brier by approximately 4.3\% and increased top-decile recall by 7.3 percentage points. At the same checklist-screening fraction, more positive reports entered the priority set, demonstrating its value for ranking monitoring records. Model rankings can identify reported locations worth revisiting, after which repeated surveys and detection models can compare occupancy persistence probabilities across elevation and greenness conditions. The following stratified analysis examines which environments benefit from improved report identification.
\inctab{spatial_valid}
\begin{figure}[htbp]\centering
\incfig{fig22_conventional_spatial_rewrite}
\caption{Conventional reporting models on 30 shared geographic test folds. Points are fold means and bars are between-fold SD, not confidence intervals. Dashed lines mark m4. The best mean overall discrimination, positive-report ranking and probability error are not all supplied by the same configuration.}\label{fig:validbench}
\end{figure}

The changes were selective (Figure~\ref{fig:validbench}). Adding terrain to m4 increased AUC from 0.800 to 0.822, while AP changed only from 0.253 to 0.255 and Brier error increased slightly. Terrain improved the boosting configurations but reduced the MLP's performance. The additive m3 and interaction m4 GLMs had almost identical spatial results despite the latter's better sample-fit AIC. Thus an environmental term can help describe the sampled pattern without providing a corresponding improvement in held-out report prediction. Different feature sets also prevent attributing tree-model gains solely to model architecture. For monitoring, terrain inputs should be judged by whether they improve positive-report coverage in target environments. Higher AUC with nearly unchanged AP indicates that stronger overall discrimination need not supply more useful revisit clues under limited screening capacity. Similar m3 and m4 performance also shows that an interaction describing rebound gradients need not improve site prioritization.

\subsection{\ifzh 低地报告识别的改善为占域监测提供复访线索\else Better identification of lowland reports guides revisits for occupancy monitoring\fi}
\par
\ifzh
低地是检验报告回升与占域持续性关系的重要环境。其报告率只有 3.3\%，却包含完整样本中 60.7\% 的正报告；同时，低海拔低、高绿度地点的模型估计的占域持续概率分别为 24.2\% 和 71.3\%。低报告率因而掩盖了低地内部占域持续概率的差异。新增十折实验显示，随机森林将低海拔 AP 从 m4 的 0.143 提高至 0.201，中海拔从 0.221 提高至 0.326，高海拔则由 0.540 变为 0.538（表~\ref{tab:envmonitor}；图~\ref{fig:envmonitor}）。预测改善主要补充了低、中海拔稀少报告的识别，能够为从既有记录中选择复访候选地点提供依据。较高占域持续概率集中于山地和较绿地点，而较大预测收益集中于低、中海拔，两者为监测提供互补信息。
\else
Lowlands are important for testing the connection between reporting rebound and occupancy persistence. Their reporting rate was only 3.3\%, yet they contained 60.7\% of complete-sample positive reports, while modeled persistence was 24.2\% and 71.3\% in low- and high-greenness lowland cells. Low aggregate reporting therefore conceals differences in occupancy persistence probabilities within lowlands. In the new ten-fold experiment, random forest increased low-elevation AP from m4's 0.143 to 0.201 and middle-elevation AP from 0.221 to 0.326; high-elevation AP changed from 0.540 to 0.538 (Table~\ref{tab:envmonitor}; Figure~\ref{fig:envmonitor}). Prediction gains chiefly improve identification of sparse low- and middle-elevation reports, providing a basis for selecting candidate revisit locations from existing records. Higher modeled persistence concentrates in uplands and greener sites, whereas larger predictive gains concentrate at low and middle elevations; these patterns supply complementary monitoring information.
\fi
\inctab{environment_monitoring}
\par
\ifzh
全局前 10\% 排序中，随机森林覆盖低海拔正报告的 34.6\%，高于 m4 的 26.4\%，差值为 8.1 个百分点，空间区块自助区间为 2.8--12.0。识别出更多低地报告，有助于扩大后续重复调查的候选范围，进而比较不同绿度地点的再次记录比例，并在考虑不完全探测后估计占域持续概率。报告识别与持续性估计需要分步完成。这些清单排序结果没有估计复访后的占域状态及其转移。高海拔正报告覆盖同时由 78.3\% 降至 72.3\%，差值为 $-6.0$ 个百分点，区间为 $[-17.2,-0.9]$。模型更换因此提高了低地报告识别，却减弱了对山地报告的覆盖；保留固定山地重复点，能够避免这一排序变化削弱对高持续性环境的长期跟踪。
\else
Under global top-decile ranking, random forest captured 34.6\% of low-elevation positives versus 26.4\% for m4, an 8.1-point gain with spatial block bootstrap interval 2.8--12.0. Identifying additional lowland reports expands the candidate pool for repeated surveys that can compare repeat-report proportions across greenness conditions and estimate occupancy persistence after accounting for imperfect detection. Report identification and persistence estimation require separate steps. Checklist ranking does not measure occupancy after a revisit. High-elevation positive coverage simultaneously declined from 78.3\% to 72.3\%, a $-6.0$-point contrast with interval $[-17.2,-0.9]$. Changing models therefore improved lowland identification while reducing upland report coverage. Retaining fixed upland repeat sites would prevent this ranking shift from weakening longitudinal coverage of environments with high modeled persistence.
\fi
\begin{figure}[htbp]\centering
\incfig{fig28_environment_monitoring}
\caption{\ifzh
预测增益与监测覆盖。（a）各海拔层的 AP；（b）随机森林两种等预算方案的层内入选比例；（c）层内正报告覆盖。两方案各选 6,890 条清单。
\else
Prediction gains and monitoring coverage. Panels show (a) AP by elevation, (b) within-band selection under two equal-budget random-forest policies, and (c) within-band positive-report coverage. Each policy selects 6,890 checklists.
\fi}\label{fig:envmonitor}
\end{figure}
\par
\ifzh
等预算比较进一步说明监测目标如何改变环境覆盖。全局随机森林排序选中低、中、高海拔清单的 6.9\%、29.7\% 和 57.9\%，将筛查名额明显集中于高报告率环境。各海拔层约选 10\% 后，低地正报告覆盖由 34.6\% 提高至 45.6\%，增加 11.0 个百分点（95\% 区间 7.4--15.5）；总体覆盖由 46.7\% 降至 35.2\%，少覆盖 388 条正报告。增加低地代表性，为比较报告回升不足与占域持续性提供了更广的记录基础，但也降低了相同筛查名额下的总体命中。用于环境差异研究时，可将固定山地复访与低地分层覆盖结合，在各层内借助模型排序选择候选地点；低地内部还应保留不同绿度的对照。当前实验量化的是按海拔分层的清单覆盖，绿度分层和实际野外成本需要在具体调查设计中另行安排。这里的合并覆盖率来自一次十折实验，与前述 30 折不加权平均分别报告。
\else
The equal-budget comparison shows how the monitoring objective changes environmental coverage. Global random-forest ranking selected 6.9\%, 29.7\% and 57.9\% of low-, middle- and high-elevation checklists, concentrating screening slots in high-reporting environments. Selecting approximately 10\% within each elevation band increased lowland positive coverage from 34.6\% to 45.6\%, an 11.0-point gain (95\% interval 7.4--15.5), while reducing overall coverage from 46.7\% to 35.2\%, or 388 fewer positive reports. Wider lowland representation supplies a broader record base for comparing incomplete reporting rebound with occupancy persistence, but reduces aggregate hits under the same screening budget. Environmental monitoring can combine fixed upland revisits with lowland coverage and use within-stratum ranking to identify candidate sites, retaining contrasting greenness conditions within lowlands. The experiment quantifies checklist coverage under elevation stratification; greenness allocation and actual field costs require a separate survey design. These pooled rates come from a single ten-fold experiment and are reported separately from the earlier unweighted 30-fold averages.
\fi
\FloatBarrier


\subsection{Spatial validation defines the geographic scope of revisit guidance}
Random validation yielded larger average performance values for the illustrated conventional configurations (Figure~\ref{fig:protocolvalid}). Random forest's AUC/AP were 0.910/0.442 under random splitting and 0.842/0.286 under spatial blocking; m4 changed from 0.820/0.299 to 0.800/0.253. The flexible model's advantage was therefore smaller in the geographic holdout. Random forest's AP advantage narrowed from 0.143 under random splitting to 0.033 under spatial splitting. This contraction weakens the basis for exporting high random-validation scores to new locations. Together with the preceding lowland analysis, the remaining spatial-holdout gains support using models to identify revisit candidates within the observed environmental domain. However, aggregate spatial performance does not establish equal gains in every lowland greenness cell. Site selection should preserve elevation and greenness coverage before applying a ranking tool assessed under geographic holdout, so that easily predicted locations do not replace the environmental contrasts needed to estimate occupancy persistence. Training fractions and sample composition also differ between protocols; the entire score decline cannot be attributed to geographic distance.
\begin{figure}[htbp]\centering
\incfig{fig23_protocols_rewrite}
\caption{Conventional-model means under random and spatial validation. Random validation has 15 evaluations and spatial validation 30. Lines connect protocol summaries, not paired treatment effects. The protocols differ in training fraction as well as geographic composition.}\label{fig:protocolvalid}
\end{figure}

\FloatBarrier
\subsection{\ifzh 时间留出支持后期复访候选记录的识别\else Temporal holdout supports identification of later-period revisit candidates\fi}
\label{sec:additionalvalidation}
\ifzh
随机、空间和时间验证显示出不同的难度结构（图~\ref{fig:proto}；表~\ref{tab:bench_random}）。随机划分中，训练与测试覆盖相近的地点组合，树模型能够利用细致的局地模式；空间留出后其优势缩小，说明这些局地模式的可迁移程度低于随机划分表现。在 Fiona 后时间测试中，随机森林的 AUC、AP 和 Brier 分别为 0.8901、0.3869 和 0.0391，而 m4 分别为 0.8321、0.3174 和 0.0422（表~\ref{tab:bench_temporal}）。随机森林将前十分位召回率从 53.28\% 提高到 59.80\%，同时降低概率误差，表明较早时期学习的环境与观测关系对末期报告仍有预测价值。利用较早记录得到的排序仍能在 Fiona 后的观测中集中较多正报告，可用于筛选值得重新调查的地点。与低地分层结果结合，其监测用途是扩大稀少报告的检索，再通过同一地点的重复访问积累探测历史，估计占域状态及其转移。不过，该时间测试汇总所有海拔层，尚未单独量化低地不同绿度组合的跨时期收益；监测方案应保留这些组合的固定复访，检验预测信息是否同样适用。时间测试沿用同一岛屿的观测网络，其较好表现支持报告线索的跨时期使用，对鸟类占域持续性的判断仍来自重复调查和转移模型。
\else
The random, spatial and temporal protocols expose different prediction demands (Figure~\ref{fig:proto}; Table~\ref{tab:bench_random}). Random splits retain similar location mixtures in training and testing, allowing tree models to use detailed local patterns. Their smaller spatial-holdout advantage indicates weaker transfer of those patterns across locations. In the Post-Fiona time test, random forest achieved AUC 0.8901, AP 0.3869 and Brier 0.0391, compared with 0.8321, 0.3174 and 0.0422 for m4 (Table~\ref{tab:bench_temporal}). Top-decile recall increased from 53.28\% to 59.80\% while probability error declined. Environmental and observation relationships learned in earlier periods therefore retained predictive value for later reports. Rankings learned from earlier records still concentrate positive reports in Post-Fiona observations and can help select sites worth surveying again. In conjunction with the lowland analysis, its monitoring role is to broaden searches for sparse reports before repeated visits supply detection histories for estimating occupancy states and transitions. The temporal test pools elevation bands and does not separately quantify transfer within lowland greenness cells; fixed revisits across those cells are needed to assess whether the predictive information applies equally. Because the test uses the same island observation network, it supports cross-period use of reporting clues, while inference about continued occupancy rests on repeated surveys and transition models.
\fi
\begin{figure}[htbp]\centering\incfig[\linewidth]{fig06_protocol_comparison}
\caption{\ifzh
三种验证情境中的 AUC 与 AP。空间、随机和时间测试分别强调地理区块留出、观测插值和时期迁移；神经配置的训练来源重叠见第~\ref{sec:neuraldiagnostic}节。
\else
AUC and AP under geographic holdout, observation interpolation and transfer across periods. Neural training-source overlap is analyzed in Section~\ref{sec:neuraldiagnostic}.
\fi}\label{fig:proto}
\end{figure}
\inctab{bench_random}
\inctab{bench_temporal}
\FloatBarrier
\subsection{\ifzh 区块尺度与植被指数影响预测增益的幅度\else Block size and vegetation index affect the size of predictive gains\fi}
\ifzh
区块边长从 $0.04^\circ$ 增至 $0.32^\circ$ 时，加入地形的 LightGBM 的 AUC 从 0.8555 降至 0.8351，AP 从 0.3066 降至 0.2592（图~\ref{fig:blocksize}）。m4 的 AUC 分别为 0.8002 和 0.8024，端点变化较小。相应的 AUC 优势由 0.0553 缩小至 0.0327，说明灵活模型捕获的信息中，一部分依赖较细的空间组织。在 $0.16^\circ$ 下，LightGBM 的 AUC 为 0.8413，高于 m4 的 0.8146，但 AP 为 0.2904，低于 m4 的 0.2975；总体判别改善并不保证稀少正报告在排序顶端更加集中。曲线的非单调变化反映了区块合并和折内样本组成共同改变评价难度，区块宽度应与目标监测覆盖尺度一并考虑。低地持续性调查需要在拟开展调查的区域内检验地点排序的表现。在较近地点间有效的细致地形和位置线索，跨更大区块后可能减弱。扩大调查范围时，保留各环境层的固定点可防止排序效能下降造成某一类环境持续缺少复访；这些全样本尺度结果尚未单独确定低地各绿度层的适用距离。
\else
As block width increased from $0.04^\circ$ to $0.32^\circ$, terrain-augmented LightGBM declined from AUC 0.8555 to 0.8351 and from AP 0.3066 to 0.2592 (Figure~\ref{fig:blocksize}). The corresponding m4 AUC values were 0.8002 and 0.8024. LightGBM's AUC advantage consequently narrowed from 0.0553 to 0.0327, indicating that part of its predictive information depends on finer spatial organization. At $0.16^\circ$, LightGBM retained higher AUC than m4 (0.8413 versus 0.8146) but lower AP (0.2904 versus 0.2975). Better overall discrimination thus did not consistently concentrate rare positive reports at the top of the ranking. Nonmonotonic changes reflect the joint effects of block aggregation and fold composition, making the monitoring scale relevant to model assessment. For lowland persistence surveys, this supports assessing candidate-site rankings at the intended regional scale. Terrain and location information useful among nearby sites may weaken across larger blocks. Fixed sites in each environmental stratum can preserve revisit coverage as the survey domain expands; these pooled scale tests do not separately establish transfer distances for lowland greenness cells.
\fi
\begin{figure}[htbp]\centering\incfig[\linewidth]{fig17_blocksize_sensitivity}
\caption{\ifzh
四种地理区块宽度下的 AUC 与 AP。各尺度采用六折评价，误差线为折均值的标准误；划分未设置距离缓冲。
\else
AUC and AP across four geographic block widths, with six folds at each width. Bars show standard errors across folds; no distance buffer was imposed.
\fi}\label{fig:blocksize}
\end{figure}
\ifzh
EVI2 的六折实验保持了树模型在报告排序上的优势（表~\ref{tab:bench_evi2}）。随机森林的 AP 为 0.3223，比同一实验中的 m4（0.2765）提高 0.0458，Brier 从 0.0402 降至 0.0386。加入地形的 XGBoost 取得最高常规模型 AUC（0.8576），前十分位召回率为 56.15\%。因此，树模型优势并不只出现在采用 NDVI 的分析中。另一方面，EVI2 下 MLP 的 Brier 为 0.0515，加入地形后升至 0.0583，显示增加地形输入没有改善其概率估计。两个植被指数的实验共同指向模型对环境信息的利用方式，而非输入维度本身，是预测表现的重要因素。替代绿度指数仍能支持报告排序，增强了以植被信息辅助候选地点筛查的依据；但该实验比较的是清单预测，并未重新估计 EVI2 与占域持续性的关系。因此，绿度指标对监测排序的帮助与绿度与占域持续概率的关联，需要分别由预测和转移结果支撑。
\else
The six-fold EVI2 experiment retained the tree models' ranking advantage (Table~\ref{tab:bench_evi2}). Random forest achieved AP 0.3223, exceeding m4 within the same experiment (0.2765) by 0.0458, while reducing Brier from 0.0402 to 0.0386. Terrain-augmented XGBoost had the highest conventional-model AUC (0.8576) and top-decile recall of 56.15\%. The tree-model advantage therefore also appeared with EVI2. In contrast, MLP Brier increased from 0.0515 to 0.0583 when terrain was added. Across the vegetation-index experiments, performance depended on how models used environmental information, rather than simply on the number of inputs. Retaining useful report ranking with an alternative vegetation index strengthens the basis for vegetation-informed screening of candidate locations. However, this experiment did not re-estimate the association between EVI2 and persistence. The value of greenness for screening and its association with occupancy persistence therefore rest on prediction and transition evidence, respectively.
\fi
\inctab{bench_evi2}
\FloatBarrier
\subsection{\ifzh 排序改善与概率校准呈现不同的模型优势\else Ranking and calibration favor different model properties\fi}
\ifzh
十折合并预测的 ROC 与精确率--召回率曲线进一步解释了性能差异（图~\ref{fig:curves}）。m4 的合并 AUC/AP 为 0.8088/0.2836，加入地形的 XGBoost 为 0.8640/0.3162，LightGBM 为 0.8625/0.3086。两类提升树均改善了正报告与未报告清单之间的排序，XGBoost 在正报告集中的排序区间取得更高 AP。由于正报告仅占完整样本的 4.89\%，精确率--召回率曲线比单独报告 AUC 更直接地体现高分清单中的命中密度。
\else
The pooled ten-fold ROC and precision--recall curves clarify the ranking differences (Figure~\ref{fig:curves}). Pooled AUC/AP were 0.8088/0.2836 for m4, 0.8640/0.3162 for terrain-augmented XGBoost and 0.8625/0.3086 for terrain-augmented LightGBM. Both boosting models improved ordering between positive and nonreporting checklists; XGBoost provided higher AP in the positive-report ranking. With reports comprising only 4.89\% of complete records, precision--recall curves directly describe the concentration of positives among highly ranked checklists.
\fi
\begin{figure}[htbp]\centering\incfig[\linewidth]{fig08_roc_pr_curves}
\caption{\ifzh
十折预测合并后的 ROC 与精确率--召回率曲线。括号内为合并 AUC 或 AP；水平虚线表示正报告比例。
\else
ROC and precision--recall curves pooled across the ten-fold prediction export. Parentheses give pooled AUC or AP; the horizontal dashed line marks report prevalence.
\fi}\label{fig:curves}
\end{figure}
\ifzh
概率校准呈现不同次序（图~\ref{fig:calib}）。m4 的期望校准误差为 0.0032，低于样条 GLM 的 0.0050、地形 XGBoost 的 0.0093 和地形 LightGBM 的 0.0181。XGBoost 的合并 Brier 为 0.0389，低于 m4 的 0.0402，但其分组平均概率与实际频率的偏离更大。该组合说明，树模型通过更强的个体区分能力降低总体平方误差，而简单 GLM 更接近分组频率尺度。用于优先排列调查清单时，应重视 AP 和前十分位召回；用于汇总预期报告数量时，则需同时关注概率校准。排序帮助找到低地稀少报告，校准则影响各环境层预期报告水平的估计。若直接把偏高的预测概率当作回升幅度，可能夸大该层的变化；目前的总体校准结果需要在环境层内继续核查，才能用于比较低地与山地的预期报告数量。
\else
Calibration produced a different ordering (Figure~\ref{fig:calib}). Expected calibration error was 0.0032 for m4, versus 0.0050 for the spline GLM, 0.0093 for terrain-augmented XGBoost and 0.0181 for terrain-augmented LightGBM. XGBoost nevertheless had lower pooled Brier than m4 (0.0389 versus 0.0402). Its stronger discrimination reduced overall squared error while grouped predicted probabilities deviated more from observed frequencies. AP and top-decile recall consequently address prioritization, whereas probability calibration also matters when aggregating expected report counts. Ranking helps identify sparse lowland reports, while calibration affects estimated reporting levels within environmental strata. Upward-biased probabilities could overstate apparent rebound if interpreted directly as change. The current pooled calibration results require within-stratum assessment before comparing expected report counts between lowlands and uplands.
\fi
\begin{figure}[htbp]\centering\incfig[\linewidth]{fig09_calibration}
\caption{\ifzh
概率可靠性与校准误差。可靠性曲线采用 12 个等频组，期望校准误差采用 15 个等频组；对角线表示预测概率与观测频率一致。
\else
Probability reliability and calibration error. Reliability curves use 12 equal-count bins and expected calibration error uses 15; the diagonal denotes agreement between predicted probability and observed frequency.
\fi}\label{fig:calib}
\end{figure}
\FloatBarrier

\subsection{\ifzh 训练来源重叠解释了神经配置的偏高评价\else Training-source overlap compromises neural performance estimates\fi}
\label{sec:neuraldiagnostic}
\ifzh
重建三次空间划分后，完整记录测试样本与对应全部记录神经训练来源的重叠比例分别为 \rwOverlapOne\%、\rwOverlapTwo\% 和 \rwOverlapThree\%（图~\ref{fig:overlap}）。其来源是两张表分别进行样本量均衡分折，使同一地理区块获得不同折编号；随后按编号筛选，部分测试记录被纳入训练。HERON-plus 未设置内部早停留出，因而直接使用这些重叠记录拟合。原 HERON 的内部留出会进一步筛选训练来源，故图中比例表示其训练来源暴露率。该诊断确定了跨表划分对应关系是评价偏差的具体来源。
\else
Across three reconstructed spatial partitions, \rwOverlapOne\%, \rwOverlapTwo\% and \rwOverlapThree\% of complete-record test observations appeared in the corresponding all-record neural training sources (Figure~\ref{fig:overlap}). Independently balancing folds in the two tables assigned different fold identifiers to the same geographic blocks. Subsequent filtering by identifier admitted test records to training. HERON-plus, which omitted internal early-stopping holdout, fitted these overlapping records directly. Original HERON further filtered its training source internally, so the displayed percentage measures source exposure for that configuration. The diagnostic identifies cross-table partition alignment as a concrete source of evaluation bias.
\fi
\begin{figure}[htbp]\centering\incfig[\linewidth]{fig24_training_overlap_rewrite}
\caption{\ifzh
三次空间划分的逐折训练来源重叠比例。每条线代表一次划分，横轴为测试折；正文的汇总比例按每次划分全部 68,976 条测试记录加权。
\else
Fold-specific training-source overlap across three spatial partitions. Each line represents one partition and the horizontal axis identifies test folds. Text summaries weight all 68,976 test records within each partition.
\fi}\label{fig:overlap}
\end{figure}
\ifzh
全模型比较中，HERON-plus 的平均 AUC/AP 为 0.8952/0.3916（图~\ref{fig:bench}；表~\ref{tab:bench}），在共同 20 折中分别有 19 折和 20 折高于 m4（图~\ref{fig:paired}）。十折直接比较中，其 AUC/AP 从 HERON 的 0.8406/0.2686 升至 0.8831/0.3624，前十分位召回率从 48.97\% 升至 59.48\%（表~\ref{tab:bench_headtohead}）。这种跨折一致的高分与广泛训练暴露同时存在，表明配对测试集并未消除训练阶段的信息共享。由此，神经配置的高分应解释为包含测试暴露的预测表现；本文对地理推广能力的判断采用训练与测试分离的常规模型结果。这种信息共享削弱了神经高分用于新增监测地点的依据，也不能加强高得分环境具有较高占域持续概率的解释。对低地复访，应依据已分离训练与测试的环境分层结果选择候选记录，使监测覆盖的改善来自留出报告识别能力。
\else
In the all-model comparison, HERON-plus had mean AUC/AP 0.8952/0.3916 (Figure~\ref{fig:bench}; Table~\ref{tab:bench}), exceeding m4 in 19 and 20 of the 20 shared folds, respectively (Figure~\ref{fig:paired}). In the ten-fold head-to-head comparison, its AUC/AP rose from HERON's 0.8406/0.2686 to 0.8831/0.3624, with top-decile recall increasing from 48.97\% to 59.48\% (Table~\ref{tab:bench_headtohead}). Consistently high scores coexisted with extensive training exposure. Pairing test sets did not remove information sharing during fitting. Neural scores describe prediction with test exposure, and the geographic-transfer assessment therefore rests on the conventional models with separated training and testing. This information sharing weakens the use of neural scores to justify new monitoring locations and cannot strengthen a persistence interpretation for high-scoring environments. Lowland revisit candidates should instead be informed by the separated-training stratified results, so that coverage gains rest on identification of held-out reports.
\fi
\begin{figure}[htbp]\centering\incfig[\linewidth]{fig05_main_benchmark_spatial}
\caption{\ifzh
共同 20 折上的全模型性能。误差线为折间标准误；神经配置的评价包含图~\ref{fig:overlap}所示的训练来源重叠。
\else
All-model performance on 20 common folds. Bars show standard errors across folds; neural evaluations include the training-source overlap in Figure~\ref{fig:overlap}.
\fi}\label{fig:bench}
\end{figure}
\inctab{bench}
\begin{figure}[htbp]\centering\incfig[\linewidth]{fig07_paired_folds}
\caption{\ifzh
HERON-plus 与 m4 的逐折配对结果。点表示共同测试折，虚线表示两模型分数相同；神经训练来源与测试样本存在重叠。
\else
Paired fold results for HERON-plus and m4. Points represent shared test folds and the dashed line denotes equal scores; neural training sources overlap test observations.
\fi}\label{fig:paired}
\end{figure}
\inctab{bench_headtohead}
\ifzh
多指标比较将这种差异分解为判别、概率误差和高分清单覆盖（图~\ref{fig:heat}）。常规树模型的改善集中在排序及部分概率误差指标，MLP 的增加维度配置则在多个指标上下降，与主基准中的 Brier 增加一致。F1、MCC、平衡准确率和召回率共同受同一组预测与阈值影响，因此这些列反映同一预测结果的不同使用方式。对监测决策而言，指标图的主要作用是识别模型在哪类任务上有收益，以及收益是否伴随概率质量损失。较高召回率可增加复访候选记录，较好的概率质量可支持预期报告量估计；两者都需要与环境分层结合，才能判断低地覆盖是否增加、山地跟踪是否减弱。多项指标一致改善不会增加地点状态转移的观测信息。
\else
The metric comparison separates discrimination, probability error and coverage of highly ranked checklists (Figure~\ref{fig:heat}). Conventional tree-model gains concentrate in ranking and selected probability-error measures, whereas the higher-dimensional MLP declines on multiple measures, consistent with its increased Brier score. F1, MCC, balanced accuracy and recall depend on the same predictions and threshold rule; they describe different uses of those predictions. For monitoring, the display identifies which tasks benefit and whether ranking gains accompany losses in probability quality. Higher recall can expand candidate records and better probability quality can support expected-report estimates. Both must be combined with environmental stratification to determine whether lowland coverage improves or upland tracking weakens. Agreement across metrics does not add observations of site-state transitions.
\fi
\begin{figure}[htbp]\centering\incfig[\linewidth]{fig18_metric_winloss_heatmap}
\caption{\ifzh
相对 m4 的 16 项预测指标差异。颜色表示按逐折差值标准差缩放的改善幅度，正负号表示改善方向；图中不进行显著性分类。
\else
Differences from m4 across 16 prediction metrics. Colors scale gains by the standard deviation of fold differences; signs indicate the direction of improvement. No significance classification is applied.
\fi}\label{fig:heat}
\end{figure}
\FloatBarrier
\subsection{\ifzh 消融结果揭示空间输入与树成员的主要作用\else Ablations highlight spatial inputs and tree members\fi}
\ifzh
神经分支消融中，空间 Fourier 特征产生最大的 AUC 增量（+0.0347），加入地形次之（+0.0204）；去除分组损失或空间风险惩罚的 AUC 变化均为 $-0.0002$（图~\ref{fig:abl}；表~\ref{tab:ablation}）。将分组损失权重增至五倍反而使 AP 下降 0.0124。该组合表明，当前拟合分数主要响应空间表达能力和环境输入，分组约束没有带来相应的排序改善。去除单调努力约束使 AUC/AP 增加 0.0101/0.0078，说明约束压缩了模型拟合空间；其价值在于规定响应形状，而非提高当前评价分数。Platt 校准保留 AUC/AP，却将 Brier 从未校准的 0.03953 变为 0.03971，表明该校准步骤没有改善此配置的总体概率误差。上述比较描述含训练暴露流程的内部响应，空间特征亦可能利用这种暴露增强拟合。对环境解释而言，空间表达带来的高分不能据此归为更强的植被或海拔效应；分组约束的小幅变化也未增加占域持续性的直接证据。该消融因此将监测用途定位在候选记录的拟合排序，而占域持续性的判断仍需重复调查中的状态转移。
\else
Spatial Fourier features produced the largest neural-branch AUC increase (+0.0347), followed by terrain (+0.0204). Removing the grouped loss or spatial-risk penalty changed AUC by only $-0.0002$ in each case (Figure~\ref{fig:abl}; Table~\ref{tab:ablation}). Increasing the grouped-loss weight fivefold reduced AP by 0.0124. Scores therefore responded mainly to spatial representation and environmental inputs, while grouped constraints supplied little ranking improvement. Removing monotone effort constraints increased AUC/AP by 0.0101/0.0078, showing the fitting restriction imposed by that structural choice. Platt calibration preserved AUC/AP but changed Brier from 0.03953 without calibration to 0.03971 with it, providing no improvement in overall probability error. These are internal responses of the exposed-training workflow; spatial features may also exploit that exposure. For environmental interpretation, gains from spatial representation cannot be assigned to stronger vegetation or elevation effects, and small grouped-loss responses add no direct persistence evidence. The ablation informs fitted record ranking; occupancy persistence still requires transition information from repeated surveys.
\fi
\begin{figure}[htbp]\centering\incfig[\linewidth]{fig10_ablation}
\caption{\ifzh
神经分支的 20 折组件比较。差值相对于基础神经分支，误差线为配对差值的标准误；各行展示该训练流程对组件改变的响应。
\else
Twenty-fold neural-component comparisons. Deltas are relative to the base neural branch and bars show standard errors of paired differences; rows describe responses to component changes within this training workflow.
\fi}\label{fig:abl}
\end{figure}
\inctab{ablation}
\ifzh
六折堆叠消融显示，去除树成员使 AUC 下降 0.0311、AP 下降 0.0271，而去除神经成员仅分别下降 0.0034 和 0.0046（表~\ref{tab:ablation_stack}）。因此，该组合的分数主要由树成员维持，神经成员的增量较小。去除 GLM 后 AUC 增加 0.0010，但 AP 下降 0.0067、Brier 增加 0.00036，说明 GLM 对总体判别的贡献较小，却有助于正报告排序和概率误差。堆叠的作用表现为成员间的互补，而非所有成员在所有指标上均提供同方向收益。其对监测的启示是优先检验较简单树模型能否完成低地候选记录筛查，再评价组合是否增加特定环境的覆盖。这里的成员贡献没有按环境分解，不能据此推断神经分支识别出了常规模型遗漏的持续占域网格。
\else
In the six-fold stack ablation, removing tree members reduced AUC by 0.0311 and AP by 0.0271, whereas removing the neural member reduced them by only 0.0034 and 0.0046 (Table~\ref{tab:ablation_stack}). Tree members therefore maintained most of the combination's score, with a smaller neural increment. Removing the GLM increased AUC by 0.0010 but reduced AP by 0.0067 and increased Brier by 0.00036. The GLM contributed little to overall discrimination while supporting positive-report ranking and probability error. Stacking consequently reflected complementary contributions across metrics. For monitoring, this favors testing whether simpler tree models suffice for lowland candidate screening before assessing whether a combination adds coverage in particular environments. These member contributions were not decomposed by environment and do not establish that the neural branch located additional grids with persistent occupancy.
\fi
\inctab{ablation_stack}
\ifzh
将训练来源从 68,976 条完整记录扩展到 89,682 条全部记录时，六折堆叠的 AUC/AP 从 0.8623/0.3343 升至 0.8658/0.3398（图~\ref{fig:miss}）。其 AP 增量为 0.0055，明显小于 20 折神经分支比较的 0.0217，说明增加记录在单分支和组合模型中的贡献并不相同。新增记录的正报告比例较高（8.71\% 对 4.89\%），扩展训练同时改变了报告组成和跨表训练暴露。因而该结果反映的是数据组成与划分关系共同变化后的分数响应，不能把增量全部归给额外缺失记录携带的环境信息。结合高海拔 NDVI 缺失率达 49.6\% 的结果，保留缺失记录具有改善山地代表性的明确数据价值；但该价值尚未被这里的分数增量单独量化。环境监测应先保证缺失记录对应地点仍被纳入覆盖，并补充可比较的植被信息，再检验其是否改变山地报告回升和占域持续概率的估计。
\else
Expanding the training source from 68,976 complete records to all 89,682 records increased six-fold stack AUC/AP from 0.8623/0.3343 to 0.8658/0.3398 (Figure~\ref{fig:miss}). The AP increase of 0.0055 was smaller than the 0.0217 contrast in the 20-fold neural-branch comparison, showing that additional records affected the branch and the combined predictor differently. Added records had a higher report rate (8.71\% versus 4.89\%), and expansion changed both report composition and cross-table training exposure. The score response therefore combines changes in data composition and partition alignment, rather than isolating environmental information supplied by the additional records. Given 49.6\% NDVI missingness at high elevation, retaining missing records has a clear representational value for uplands, although these score increments do not isolate that benefit. Monitoring should preserve coverage of locations associated with missing records and recover comparable vegetation information before testing whether they change estimated upland reporting rebound and occupancy persistence.
\fi
\begin{figure}[htbp]\centering\incfig[\linewidth]{fig20_missingness_gain}
\caption{\ifzh
六折堆叠中全部记录与完整记录训练的比较。线连接同一折的结果；扩展数据同时改变样本组成与训练来源对应关系。
\else
All-record versus complete-record training in the six-fold stack. Lines connect results from the same fold; expansion changes both sample composition and training-source alignment.
\fi}\label{fig:miss}
\end{figure}
\FloatBarrier

\subsection{\ifzh 努力响应、环境曲线与海拔质心揭示模型的拟合结构\else Effort responses, environmental curves and centroids reveal fitted structure\fi}
\label{sec:responses}
\ifzh
努力分支得分随调查时长从 5 分钟增至 240 分钟而由 0.286 增至 0.588，随距离从 0 增至 5~km 而由 0.316 增至 0.533（图~\ref{fig:eff}）。这使同一环境条件下投入较多观察努力的清单获得较高报告预测，显示努力标准化在时期比较中的必要性。曲线方向由单调约束规定，拟合决定其斜率和弯曲程度；因此，该图用于解释报告得分如何随努力量调整。若低地与山地或前后时期的访问时长不同，未经努力统一的报告差异就可能被放大或减弱。复访应保持可比较的调查时长和距离，并保留未报告访问，并结合探测模型，才能检验报告增加是否伴随较高的占域持续概率。
\else
The effort-branch score increased from 0.286 to 0.588 as duration rose from 5 to 240 minutes, and from 0.316 to 0.533 as distance rose from 0 to 5~km (Figure~\ref{fig:eff}). Under otherwise fixed conditions, more observation effort therefore yields a higher report prediction. This illustrates the role of effort standardization in period comparisons. Monotonic constraints set the response direction, while fitting determines slope and curvature; the curves explain how effort adjusts the reporting score. Differences in visit duration across elevations or periods could amplify or weaken unstandardized reporting contrasts. Comparable duration and distance, together with retained nonreporting visits, allow repeat surveys and detection models to test whether increased reporting accompanies higher occupancy persistence.
\fi
\begin{figure}[htbp]\centering\incfig[\linewidth]{fig11_detection_effort_curve}
\caption{\ifzh
神经模型的努力分支响应。其他输入固定，单调约束规定得分随时长和距离非递减。
\else
Effort-branch responses with other inputs fixed. Monotonic constraints require nondecreasing scores with duration and distance.
\fi}\label{fig:eff}
\end{figure}
\ifzh
海拔响应呈明显非线性。神经环境分支在较高海拔接近饱和，GLM 的线性 logit 响应则持续增加（图~\ref{fig:pd}）。NDVI 分支呈单峰形状，三个时期的报告响应峰值分别位于约 0.53、0.47 和 0.51，表明灵活拟合将最高绿度与最高得分分离。该形状与 GLM 的平均正关联共同说明，整体梯度与局部曲率是不同的描述层次；神经单峰响应提供了检查高绿度范围样本覆盖及条件组合的具体线索。对监测布点，这种模型间差异支持同时保留中、高绿度地点，以检验报告曲率是否在重复调查中出现。当前神经训练暴露及分支得分定义削弱了把峰值解释为最适持续性绿度的依据，因此不宜据此减少高绿度地点的复访。
\else
Elevation responses were nonlinear. The neural environmental branch approached saturation at higher elevations, whereas the GLM's linear-logit response continued rising (Figure~\ref{fig:pd}). The NDVI response was unimodal, with reporting-response peaks near 0.53, 0.47 and 0.51 across the three periods. Flexible fitting thus separated the greenest conditions from the highest fitted score. Together with the GLM's average positive association, these curves distinguish an overall gradient from local curvature and identify high-greenness sample coverage and covariate combinations as relevant checks on the fitted peak. For survey placement, this model disagreement supports retaining both intermediate- and high-greenness sites to test whether the reporting curvature appears in repeat observations. Training exposure and the branch-score definition weaken an interpretation of the peak as optimal greenness for persistence, providing no basis to reduce revisits in the greenest environments.
\fi
\begin{figure}[htbp]\centering\incfig[\linewidth]{fig12_psi_partial_dependence}
\caption{\ifzh
GLM 报告曲线与神经模型的报告、环境分支响应。实线表示神经拟合，虚线表示 GLM；分支得分采用条件输入绘制。
\else
GLM reporting curves and neural reporting and environmental-branch responses. Solid lines denote neural fits and dashed lines denote the GLM; branch scores use fixed conditional inputs.
\fi}\label{fig:pd}
\end{figure}
\ifzh
这种曲线形状差异解释了绝对海拔质心的差别（图~\ref{fig:cent}）。GLM 的三个时期质心为 939.7、968.2 和 953.5~m，神经报告曲线为 797.8、838.7 和 817.4~m。高海拔饱和降低了神经曲线对预测网格上端的相对权重，因此其绝对质心更低。两个模型的海拔质心均先升高，随后部分回落。GLM 相对初期的位移为 +28.5 和 +13.9~m，神经模型为 +40.9 和 +19.6~m。这一共同次序概括了报告权重的时期重分配，且位移幅度对响应函数形状敏感。时期比较因此需要同时覆盖低地和山地。只跟踪高海拔高报告地点，会遗漏后期报告权重向较低海拔回落所涉及的变化。两条曲线来自共享观测，不能把共同次序当作独立验证；同地点的探测历史与状态转移模型才有助于检验这种重分配是否伴随低地占域状态的维持。
\else
Response shape explains the different absolute elevation centroids (Figure~\ref{fig:cent}). GLM centroids were 939.7, 968.2 and 953.5~m, versus 797.8, 838.7 and 817.4~m for the neural reporting curve. High-elevation saturation reduced the neural curve's relative weight at the upper end of the prediction grid, lowering its absolute centroid. Both models showed an upward centroid shift between hurricanes followed by a partial return. Shifts relative to the initial period were +28.5 and +13.9~m for the GLM, and +40.9 and +19.6~m for the neural model. This shared ordering summarizes period-specific redistribution of fitted reporting weights, while its magnitude depends on response shape. Period comparisons therefore require both lowland and upland coverage. Tracking only high-reporting upland locations could miss record changes associated with the later downward reweighting. The curves share observations and their common ordering is not independent validation; same-site detection histories and transition models can test whether the redistribution accompanies persistent lowland occupancy.
\fi
\begin{figure}[htbp]\centering\incfig[0.74\linewidth]{fig13_elevational_centroid}
\caption{\ifzh
0--1300~m 预测网格上的概率加权海拔质心。标注为各模型相对自身初期基线的位移。
\else
Probability-weighted elevation centroids on the 0--1300~m prediction grid. Annotations show shifts relative to each model's own initial baseline.
\fi}\label{fig:cent}
\end{figure}
\FloatBarrier
\subsection{\ifzh 空间异质性与变量贡献对应不同监测信息\else Spatial heterogeneity and variable contributions inform monitoring\fi}
\ifzh
复现的原占域模型中，在平均绿度下局地灭绝预测从近海平面的 56.7\% 降至约 1300~m 的 0.9\%（图~\ref{fig:ext}）。高海拔端区间较宽，表明极端环境条件下的幅度估计受样本信息限制；观测排列校正后的海拔系数仍为负值，与这一梯度方向一致。神经环境分支补数也随海拔下降，但它在数学上只是低环境得分的表达。两个面板分别展示地点转移估计与条件得分变化，持续性解释来自动态占域模型的转移分量。海拔梯度因而为固定山地复访提供了过程依据，而低地较高的转移风险使低地绿度对照更有必要。结合校正后低地高绿度地点较高的占域持续概率，应检验哪些低地条件与较高的占域持续概率相关，避免把整个低地统一解释为高流失环境。
\else
In the reproduced original occupancy model, predicted local extinction at mean greenness declined from 56.7\% near sea level to 0.9\% at approximately 1300~m (Figure~\ref{fig:ext}). The broad upper-elevation interval indicates limited information about the magnitude at environmental extremes. The negative elevation coefficient after observation-order correction retained the same gradient direction. The complement of the neural environmental branch also declined with elevation, but mathematically expresses a low environmental score. The panels thus show a site-transition estimate and a conditional score response; the persistence interpretation derives from the dynamic model's transition component. The elevation gradient supplies a process-based rationale for fixed upland revisits, while higher modeled lowland transition risk motivates lowland greenness contrasts. Together with higher corrected persistence at greener lowland sites, it supports testing which lowland conditions are associated with higher occupancy persistence probabilities rather than treating all lowlands as high-loss environments.
\fi
\begin{figure}[htbp]\centering\incfig[\linewidth]{fig14_extinction_vs_elevation}
\caption{\ifzh
海拔梯度上的两种模型输出。（a）复现的原动态占域模型局地灭绝概率及区间；（b）一减神经环境分支得分。
\else
Two model outputs along elevation. (a) Reproduced original dynamic-model local extinction probability and interval; (b) one minus the neural environmental-branch score.
\fi}\label{fig:ext}
\end{figure}
\ifzh
地点地图把上述梯度转换为监测布点信息（图~\ref{fig:surf}）。占域网格的转移预测与 3,727 个采样地点的神经得分都呈现空间异质性，而地点分布并不均匀。地图适合识别已有观测范围内的高低得分对照和采样空缺，支持同时覆盖山地与低地的分层布点。神经地图采用地点跨时期平均 NDVI，因此表达的是固定地点环境条件下的时期得分，植被随时间恢复的过程应由同步遥感与重复鸟类调查另行测量。可根据地图选择不同环境中的复访地点。在山地较高占域持续概率的网格中安排重复调查，低地不同绿度地点则用于检验报告回升不足是否与较高的占域持续概率并存，观测空缺则提示需要补充覆盖。具体候选点仍需核查可达性和调查条件，地图得分本身不构成地点恢复的证据。
\else
Site maps translate these gradients into monitoring-placement information (Figure~\ref{fig:surf}). Transition predictions at occupancy grids and neural scores at 3,727 sampled sites both vary spatially, while observation locations are unevenly distributed. The maps identify score contrasts and sampling gaps within the observed domain, supporting strata that include both upland and lowland settings. The neural map uses cross-period mean site NDVI and therefore expresses period scores under fixed site conditions; time-varying vegetation recovery requires concurrent remote-sensing and repeated bird observations. The maps can guide selection of revisit sites across environments. Repeat surveys in upland grids with high modeled occupancy persistence can track occupancy transitions, while lowland greenness contrasts can test whether incomplete reporting rebound coexists with high occupancy persistence, and observation gaps identify coverage needs. Candidate locations still require assessment of access and survey conditions; mapped scores themselves do not demonstrate site recovery.
\fi
\begin{figure}[htbp]\centering\incfig[\linewidth]{fig15_spatial_surfaces}
\caption{\ifzh
地点尺度的空间异质性。上图为复现的原动态模型局地灭绝预测，下图为 Fiona 后采样地点的神经环境分支得分。
\else
Site-level spatial heterogeneity. The upper panel shows reproduced original dynamic-model local extinction predictions. The lower panel shows neural environmental-branch scores at sampled sites in the Post-Fiona period.
\fi}\label{fig:surf}
\end{figure}
\ifzh
置换重要性揭示两种模型依赖的信息不同（图~\ref{fig:imp}）。神经环境分支中，打乱海拔使 AP 下降 0.0809，远大于打乱绿度的 0.0075，表明该分支主要沿海拔组织拟合得分。地形 LightGBM 的留出折中，时长贡献最大（AP 下降 0.0898），随后是经度（0.0354）、大尺度平均坡度（0.0246）和纬度（0.0212）。这些结果将树模型的预测收益具体联系到观察努力、地理位置及多尺度地形的组合，也解释了增加环境维度为何不能自动改善所有模型。监测时统一努力记录并保持空间覆盖，有助于区分观察机会与地点条件所贡献的报告差异。
\else
Permutation importance exposed different information dependencies (Figure~\ref{fig:imp}). Permuting elevation reduced neural-branch AP by 0.0809, versus 0.0075 for greenness, identifying elevation as the branch's main organizing gradient. In the held-out terrain-LightGBM fold, duration contributed most (AP reduction 0.0898), followed by longitude (0.0354), broad-scale mean slope (0.0246) and latitude (0.0212). Tree-model predictions therefore depended on a combination of observation effort, geographic position and multiscale terrain. Recording effort consistently while preserving spatial coverage helps separate reporting differences associated with observation opportunity from those associated with site conditions.
\fi
\begin{figure}[htbp]\centering\incfig[\linewidth]{fig19_permutation_importance}
\caption{\ifzh
置换变量造成的 AP 下降，误差线为置换重复的标准差。神经面板使用拟合来源样本，LightGBM 面板使用一个地理留出折，分别描述各模型的信息依赖。
\else
AP reductions after feature permutation, with standard deviations across permutations. The neural panel uses fitting-source observations and the LightGBM panel uses one geographic holdout; each describes its model's information dependence.
\fi}\label{fig:imp}
\end{figure}
\FloatBarrier

